%% kapa: exam=1 round=33 year=2022 period=1 form=A subject=real_estate subject_ko=부동산학원론 seq=11 number=91
%% source: Q-Net 감정평가사 기출문제 (cst00309) artlSeq=5209997 file=raw/1008/5209997/2022년도 제33회 감정평가사 1차시험 1교시(민법,경제학,부동산학).pdf
%% conversion: kordoc 4.24.51 -> tools/gen1.py (auto); review=proofread-sonnet; flags=none
\begin{question}{91}
분양가상한제로 인해 발생할 수 있는 문제점과 그 보완책을 연결한 것으로 옳지 않은 것은?
\begin{choices}
  \item 분양주택의 질 하락 - 분양가상한제의 기본 건축비 현실화
  \item 분양주택 배분 문제 - 주택청약제도를 통한 분양
  \item 분양프리미엄 유발 - 분양주택의 전매제한 완화
  \item 신규주택 공급량 감소 - 공공의 저렴한 택지 공급
  \item 신규주택 공급량 감소 - 신규주택건설에 대한 금융지원
\end{choices}
\end{question}
\answer{3}
